% ============================
%   Abstract
% ============================
\section*{Abstract}
\addcontentsline{toc}{section}{Abstract}

This report describes the current state of \nidsname, a Network Intrusion
Detection System (NIDS) that combines classical signature-based rule
matching with a supervised / unsupervised machine-learning detector and
exposes both as a self-contained Security Information and Event Management
(SIEM) console. The system captures network packets in real time
(via Scapy~\cite{scapy} or a deterministic synthetic traffic generator),
evaluates each packet against a YAML-seeded rule base, scores it with a
scikit-learn~\cite{pedregosa2011scikit} classifier (Random
Forest~\cite{breiman2001random} or Isolation Forest~\cite{liu2008isolation})
and persists, for every triggered alert, a ring-buffered packet snapshot to
disk so that the analyst can replay the events surrounding the
incident---a built-in forensics layer that distinguishes this project from
the typical academic IDS prototype.

The current implementation reaches Milestone~3 of the planned roadmap: the
backend (FastAPI~\cite{ramirez2022fastapi}, async SQLAlchemy, JWT
authentication, WebSocket live streams) and the React dashboard
(Vite~\cite{vite}, TailwindCSS, Recharts) are end-to-end functional, the
hybrid detection pipeline produces alerts on synthetic attack traffic, and
a one-click cross-platform launcher abstracts the setup steps. This
document summarises the motivation, theoretical background, architecture,
implementation choices and the work remaining before the final thesis
defense.
\vspace{0.4cm}

\noindent\textbf{Keywords:} Intrusion Detection, Network Security,
Machine Learning, Random Forest, Isolation Forest, SIEM, Forensics,
FastAPI, React.
